\documentclass[a4paper,12pt]{article}
\usepackage[utf8]{inputenc}
\usepackage[portuguese]{babel}
\usepackage{graphicx}
\usepackage{amsmath}
\usepackage{enumitem}
\usepackage{natbib}
\usepackage{indentfirst}
\usepackage{hyperref}
\usepackage{fancyhdr}
\usepackage{listings}
\usepackage{xcolor}
\usepackage{acronym}
\usepackage{dirtree}

\lstset{
    language=Java,
    basicstyle=\ttfamily\small,
    keywordstyle=\color{blue},
    stringstyle=\color{red},
    commentstyle=\color{green!50!black},
    numberstyle=\tiny,
    numbers=left,
    stepnumber=1,
    numbersep=5pt,
    backgroundcolor=\color{gray!10},
    showspaces=false,
    showstringspaces=false,
    showtabs=false,
    tabsize=2,
    captionpos=b,
    breaklines=true,
    breakatwhitespace=true,
    breakautoindent=true,
    escapeinside={\%*}{*)},
    linewidth=\textwidth,
    basewidth=0.5em,
}

\pagestyle{fancy}
\fancyhf{}
\fancyfoot[R]{\thepage}
\renewcommand{\headrulewidth}{0pt}

\begin{document}

% Logo da universidade
\begin{figure}
\centering
\includegraphics{Anexos/UBI_LOGO.jpg}
\end{figure}

% Título e autores
\title{Relatório do Trabalho Prático 1 - Farmácia UBI}
\author{Afonso Esteves - 55047\\Tiago Pereira - 55019\\Felipe Dourado - 53394}
\date{\today}

\maketitle
\newpage

% Sumário
\tableofcontents

\newpage

\section*{Resumo}
\hspace{10pt}
Este relatório detalha o desenvolvimento de uma aplicação de consola em Java para a gestão de uma farmácia. O sistema, desenvolvido no âmbito da disciplina de Programação Orientada a Objetos, permite gerir clientes, produtos e encomendas. A aplicação foi estruturada de forma a aplicar os conceitos fundamentais da POO, como encapsulamento, abstração e modelagem de entidades do mundo real em classes. O sistema oferece uma interface de linha de comandos para interagir com as funcionalidades implementadas.
\addcontentsline{toc}{section}{Resumo}
\section*{Palavras-chave e \textit{Keywords}}
Java; Programação Orientada a Objetos; Sistema de Gestão; Farmácia; Consola.

\newpage

\section{Introdução}
\subsection{Âmbito, enquadramento e motivação}
\hspace{10pt}
Este documento consiste num relatório técnico que descreve o desenvolvimento de uma aplicação de consola em Java, no contexto da unidade curricular de Programação Orientada a Objetos. A aplicação simula um sistema de gestão para uma farmácia, permitindo realizar operações básicas como o registo e a gestão de clientes, o controlo de stock de produtos e o processamento de encomendas.

A principal motivação para o desenvolvimento deste projeto foi a aplicação prática dos conceitos de Programação Orientada a Objetos (POO) aprendidos, traduzindo um problema do mundo real para um modelo de software coeso e funcional. O projeto permitiu explorar a modelação de entidades, a definição de responsabilidades entre classes e a criação de uma interface de utilizador textual.

O projeto situa-se na área de Desenvolvimento de Software, com foco na Programação Orientada a Objetos em Java e no desenvolvimento de aplicações de consola.

\section{Contextualização do Projeto}

\subsection{Identificação dos Desenvolvedores}
Os autores do presente projeto são estudantes do curso de Informática Web, Móvel e na Nuvem da \ac{UBI}, matriculados na disciplina de Programação Orientada a Objetos.

\subsubsection{Desenvolvedores Envolvidos}
\begin{itemize}
    \item \textbf{Afonso Esteves (55047)}
    \item \textbf{Felipe Dourado (53394)}
    \item \textbf{Tiago Pereira (55019)}
\end{itemize}     

\subsection{Contextualização Académica}
O projeto foi desenvolvido no contexto da disciplina de Programação Orientada a Objetos, lecionada pelo Professor Paul Andrew Crocker, com o objetivo de aplicar e consolidar conhecimentos em:
\begin{itemize}
    \item Modelação de problemas com base em objetos.
    \item Implementação de classes, atributos e métodos em Java.
    \item Utilização de conceitos como encapsulamento, construtores e relacionamentos entre classes.
    \item Gestão de coleções de objetos.
    \item Desenvolvimento de uma interface de utilizador baseada em texto.
\end{itemize}

\section{Especificações Técnicas}

\subsection{Ambiente de Desenvolvimento}
\begin{itemize}
    \item \textbf{Linguagem:} Java
    \item \textbf{Ambiente de Execução:} Java Development Kit (JDK)
    \item \textbf{Ambiente de Desenvolvimento:} IntelliJ IDEA
    \item \textbf{Sistema Operacional:} macOS, Windows
\end{itemize}

\section{Utilização}
\subsection{Requisitos Prévios}
\begin{itemize}
    \item \textbf{Java:} É necessário ter o JDK (Java Development Kit) instalado no sistema.
\end{itemize}

\subsection{Como executar o programa}
\begin{itemize}
    \item \textbf{Compilação:}
    \begin{enumerate}
        \item Navegue até à diretoria `TP-01/src`.
        \item Compile os ficheiros Java com o seguinte comando:
            \begin{verbatim}
javac *.java
            \end{verbatim}
    \end{enumerate}
    \item \textbf{Execução:}
    \begin{enumerate}
        \item Após a compilação, execute a aplicação com o comando:
            \begin{verbatim}
java Main
            \end{verbatim}
    \end{enumerate}
    \item \textbf{Acesso:}
    \begin{itemize}
        \item A aplicação será executada na consola, apresentando o menu principal.
    \end{itemize}
\end{itemize}

\subsection{Uso do Programa}
\begin{itemize}
    \item Ao iniciar a aplicação, será apresentado um menu principal com as opções de gestão de clientes, produtos e encomendas.
    \item O utilizador pode navegar pelos submenus para adicionar, remover, listar ou editar as informações.
    \item A interação é feita através da introdução de números correspondentes às opções de menu e do preenchimento dos dados solicitados.
\end{itemize}

\section{Estrutura da Solução Construída}
O sistema foi arquitetado com base nos princípios da Programação Orientada a Objetos, dividindo as responsabilidades em diferentes classes que representam as entidades do domínio do problema.

\subsection{Componentes do Sistema}
\textbf{Classes de Modelo:} 
\begin{itemize}
    \item \texttt{Cliente.java}: Representa um cliente da farmácia, com nome, NIF e morada.
    \item \texttt{Produto.java}: Representa um produto, com ID, nome, preço e stock.
    \item \texttt{Encomenda.java}: Modela uma encomenda, associando um cliente a um conjunto de produtos e registando o estado da mesma.
    \item \texttt{Morada.java}: Representa a morada de um cliente.
    \item \texttt{Farmacia.java}: Classe central que gere as listas de clientes, produtos e encomendas, e contém a lógica de negócio principal.
    \item \texttt{Main.java}: Ponto de entrada da aplicação, responsável pela interface com o utilizador (menus de consola).
\end{itemize}

\subsection{Estrutura de Ficheiros do Projeto}
\dirtree{%
.1 TP-01/
.2 src/.
.3 Cliente.java.
.3 Encomenda.java.
.3 Farmacia.java.
.3 Main.java.
.3 Morada.java.
.3 Produto.java.
}

\section{Funcionalidades Implementadas}
As funcionalidades implementadas na aplicação incluem:
\begin{itemize}
    \item \textbf{Gestão de Clientes:}
    \begin{itemize}
        \item Adicionar, remover, listar e editar clientes.
        \item Cada cliente é identificado por um NIF único.
    \end{itemize}
    \item \textbf{Gestão de Produtos:}
    \begin{itemize}
        \item Adicionar, remover, listar e editar produtos.
        \item Controlo de stock (adicionar e remover unidades).
    \end{itemize}
    \item \textbf{Gestão de Encomendas:}
    \begin{itemize}
        \item Criar uma nova encomenda para um cliente, adicionando produtos do stock.
        \item Listar todas as encomendas registadas.
        \item Atualizar o estado de uma encomenda (ex: "Em processamento", "Pago", "Enviado", "Entregue").
    \end{itemize}
\end{itemize}

\section{Implementação de Funcionalidades}

\subsection{Classe Farmacia: O Coração do Sistema}
A classe `Farmacia` centraliza a lógica de negócio. Ela contém as listas de clientes, produtos e encomendas e os métodos para manipular esses dados.

\begin{minipage}{\linewidth}
\begin{lstlisting}[caption={Método para adicionar um cliente em Farmacia.java}]
public void adicionarCliente(String nome, String nif, Morada morada) {
    if (acharClienteNif(nif) != null) {
        System.out.println("Erro: Já existe um cliente com o NIF " + nif);
        return;
    }
    Cliente novoCliente = new Cliente(nome, nif, morada);
    this.clientes.add(novoCliente);
    System.out.println("Cliente adicionado com sucesso!");
}
\end{lstlisting}
\end{minipage}

\subsection{Classe Main: Interface com o Utilizador}
A classe `Main` é responsável por apresentar os menus ao utilizador e recolher o seu input, servindo de ponte para as funcionalidades da classe `Farmacia`.

\begin{minipage}{\linewidth}
\begin{lstlisting}[caption={Menu principal em Main.java}]
public static void main(String[] args) {
    int opcao;
    do {
        System.out.println("\n--- BEM-VINDO À FARMÁCIA UBI ---");
        System.out.println("1. Gestão de Clientes");
        System.out.println("2. Gestão de Produtos");
        System.out.println("3. Gestão de Encomendas");
        System.out.println("0. Sair");
        System.out.print("\nEscolha uma opção: ");

        opcao = scanner.nextInt();
        scanner.nextLine();

       // Um switch para cada opção do menu.
        switch (opcao) {
            case 1:
                menuClientes();
                break;
            case 2:
                menuProdutos();
                break;
            case 3:
                menuEncomendas();
                break;
            case 0:
                System.out.println("Saindo do programa...");
                break;
            default:
                System.out.println("Opção inválida!");
        }
    } while (opcao != 0);

    scanner.close();
}
\end{lstlisting}
\end{minipage}

\subsection{Classe Encomenda: Gestão de Estados}
A classe `Encomenda` utiliza um `enum` para gerir o estado do pedido, o que torna o código mais legível e seguro.

\begin{minipage}{\linewidth}
\begin{lstlisting}[caption={Enum EstadoEncomenda em Encomenda.java}]
public enum EstadoEncomenda {
    EM_PROCESSAMENTO("Em processamento"),
    PAGO("Pago"),
    ENVIADO("Enviado"),
    ENTREGUE("Entregue");

    private final String descricao;
    EstadoEncomenda(String descricao) {
        this.descricao = descricao;
    }
    public String getDescricao() {
        return descricao;
    }
}
\end{lstlisting}
\end{minipage}

\section{Considerações de Implementação}

\subsection{Princípios de Desenvolvimento}
\begin{itemize}
    \item \textbf{Separação de Responsabilidades:} A classe `Main` trata da interação com o utilizador, enquanto a `Farmacia` e as outras classes de modelo cuidam da lógica de negócio e dos dados.
    \item \textbf{Encapsulamento:} Os atributos das classes são privados e o acesso é feito através de métodos públicos (getters e setters), protegendo o estado interno dos objetos.
    \item \textbf{Validação de Dados:} Foram implementadas validações básicas para garantir a consistência dos dados, como a verificação de NIFs duplicados ou stocks negativos.
\end{itemize}

\subsection{Desafios Encontrados}
Durante o desenvolvimento do projeto, alguns desafios foram encontrados:
\begin{itemize}
    \item \textbf{Gestão do Input do Utilizador:} Lidar com diferentes tipos de input (texto, números) e prevenir erros de execução devido a formatos inesperados exigiu atenção no uso da classe `Scanner`.
    \item \textbf{Modelação do Relacionamento entre Classes:} Definir como as classes se relacionariam (ex: uma `Encomenda` tem um `Cliente` e vários `Produtos`) foi um passo crucial para a correta estruturação do sistema.
    \item \textbf{Consistência dos Dados:} Garantir que o stock de um produto era atualizado corretamente ao criar uma encomenda foi um ponto importante para a integridade do sistema.
\end{itemize}

\section{Conclusão}
O desenvolvimento deste projeto permitiu aplicar de forma prática e consolidada os conceitos fundamentais da Programação Orientada a Objetos. A criação de um sistema de gestão de farmácia, mesmo que simplificado, demonstrou a eficácia da modelação de entidades do mundo real em classes e objetos para a construção de software modular e de fácil manutenção.

A experiência adquirida na estruturação das classes, na definição de responsabilidades e na implementação de uma interface de utilizador textual foi de grande valor para o aprofundamento dos conhecimentos em Java e no paradigma da POO.

\clearpage

\section{Referências Bibliográficas}
\begin{itemize}
    \item Schildt, Herbert. \textit{Java: A Beginner's Guide}. 6th ed., McGraw-Hill Education, 2014.
    \item Olsson, Mikael. \textit{Java Quick Syntax Reference}. Apress, 2014.
    \item Roberts, Eric S. \textit{The Art and Science of Java}. Addison-Wesley, 2008.
\end{itemize}

\section*{Acrónimos}
\begin{acronym}[UBI]
    \acro{UBI}{Universidade da Beira Interior}
\end{acronym}

\end{document}